\section{需求分析}\label{sec:requirements}
\ChapterOwner{B}{C}
\subsection{项目背景与目标}
\Writing{B}{解释历史问答难以整理、资料来源难追溯、练习与学习状态割裂的问题。给出可检验目标，说明系统管理学习事实，花园用于呈现。}
\subsection{用户角色与使用场景}
\Writing{B}{用一次学习过程说明需求。场景图放在这里，可以画从学生提问到回找原回答的操作，或用用例图概括角色与功能。}
\subsection{功能需求与优先级}
\Writing{B}{将答疑、时间/关键词/知识点历史检索列为课程必选，资料和训练列为推荐，成长与协作列为增强。每项给编号、输入、输出及验收条件。}
\subsection{数据需求与业务规则}
\Writing{B，A 核对}{从会话归属、引用具体片段、重复提交和撤销证据等实际例子里提炼规则，说清能观察到什么。}
\subsection{非功能需求}
\Writing{B}{持久化、隔离、并发一致性、失败恢复、可复现、移动端和可访问性。性能目标与实验结果分开，不把目标当成已达成。}
\subsection{需求追踪与系统范围}
\Writing{B，C 复核}{建立需求编号、实现模块和测试编号的对应表。说明 demo/live、班级 2D SVG、资料类型和生成题型的实际范围。}
\Evidence{docs/requirements-matrix.md; docs/features-explained.md}
